% =====================================================================
%  Bölüm 8 — Kümeler, İlişkiler ve Fonksiyonlar
% =====================================================================
\chapter{Kümeler, İlişkiler ve Fonksiyonlar}\label{bol:kumeler}

\begin{bolumhedefleri}
\begin{itemize}
  \item \Mathlib{}'i projeye eklemeyi, adlandırma kurallarını ve teorem
        arama araçlarını,
  \item \leantac{ring}, \leantac{linarith}, \leantac{norm\_num},
        \leantac{positivity} ve \leantac{tauto} taktiklerini,
  \item Kümeleri \emph{yüklem} olarak görmeyi ve küme cebirinin
        önermeler mantığının bir yansıması olduğunu,
  \item \leantac{ext} taktiğiyle küme eşitliği ispatlamayı,
  \item İlişkileri, denklik ilişkilerini ve kısmi sıralamaları,
  \item Birebir, örten ve eşleme fonksiyonları; bileşke ve tersleri,
  \item Cantor teoremini ve köşegen yöntemini.
\end{itemize}
\end{bolumhedefleri}

Önceki bölümlerde her şeyi sıfırdan kurduk: bağlaçları, niceleyicileri,
doğal sayıları. Bu yaklaşım, mantığın nasıl çalıştığını anlamak için
vazgeçilmezdir; ama gerçek matematik yapmak için yavaştır. Bu bölümde
\Lean{} topluluğunun yıllardır biriktirdiği büyük matematik kütüphanesi
\Mathlib{}'i kullanmaya başlıyoruz \citep{mathlib2020}. Ardından
matematiğin üç temel dilini, yani kümeleri, ilişkileri ve fonksiyonları
mantık açısından ele alacağız.

% ---------------------------------------------------------------------
\section{Mathlib'e giriş}\label{sec:mathlib}

\Mathlib{}, \Lean{}'de yazılmış, gönüllü bir topluluk tarafından
geliştirilen ve sürekli büyüyen bir matematik kütüphanesidir. Cebirden
analize, topolojiden sayılar kuramına uzanan yüz binlerce tanım ve
teorem içerir. Kütüphaneye eklenen her teorem, \Lean{}'in çekirdeği
tarafından denetlenmiştir.

\paragraph{Kurulum}
Eşlik eden projede \Mathlib{} bağımlılığı bu bölümle birlikte açıldı.
\Mathlib{} belirli bir \Lean{} sürümüyle çalıştığı için kurulumda üç adım
gerekir: projenin \lc{lean-toolchain} dosyasını \Mathlib{}'inkiyle
eşlemek, \lc{lake update} ile bağımlılıkları indirmek ve
\lc{lake exe cache get} ile önceden derlenmiş dosyaları almak. Adımların
tam hâli projenin \lc{README.md} dosyasındadır.

\begin{dikkat}[Dikkat: Önbelleği indirmeyi unutmayın]
\Mathlib{}'i sıfırdan derlemek güçlü bir bilgisayarda bile saatler sürer.
\lc{lake exe cache get} komutu, önceden derlenmiş dosyaları birkaç
dakikada indirir. Bir \Lean{} dosyasının başında \lc{import Mathlib}
satırı varsa ve dosya açılırken bilgisayarınız uzun süre ``derliyor''
görünüyorsa, büyük olasılıkla bu adım atlanmıştır.
\end{dikkat}

\leandosya[Mathlib'e ilk bakış]{Bolum08}{mathlib_ilk}

\subsection{Mathlib'in taktikleri}

\Mathlib{}, çekirdek \Lean{}'in taktiklerine güçlü yeni araçlar ekler
(\cref{tab:mathlib-taktik}).

\begin{table}[htbp]
\centering
\renewcommand{\arraystretch}{1.45}
\rowcolors{2}{tablosatir}{white}
\begin{tabularx}{\linewidth}{>{\raggedright\arraybackslash}p{2.4cm}
                             >{\raggedright\arraybackslash}X
                             >{\raggedright\arraybackslash}p{4.2cm}}
\rowcolor{tabloust}
\tablobaslik{Taktik} & \tablobaslik{Ne yapar?} & \tablobaslik{Örnek hedef} \\
\leantac{ring}        & Değişmeli halkalarda polinom eşitliklerini kanıtlar. & $(a+b)^2 = a^2 + 2ab + b^2$ \\
\leantac{linarith}    & Hipotezlerden doğrusal eşitsizlikleri türetir. & $x < y,\; y < 3 \vdash x < 3$ \\
\leantac{norm\_num}   & Somut sayısal ifadeleri hesaplar; asallık gibi özellikleri denetler. & \lc{Nat.Prime 97} \\
\leantac{positivity}  & Bir ifadenin pozitif ya da negatif olmadığını gösterir. & $0 \leq x^2 + 1$ \\
\leantac{tauto}       & Önermeler mantığının totolojilerini (klasik olarak) kanıtlar. & $P \wedge Q \to Q \vee R$ \\
\leantac{ext}         & İki küme ya da fonksiyonun eşitliğini, elemanlarının eşitliğine indirger. & $s \cap t = t \cap s$ \\
\lc{exact?}           & Kütüphanede hedefi doğrudan kapatan bir teorem arar. & \lc{a + b = b + a} \\
\end{tabularx}
\caption{\Mathlib{}'in sık kullanılan taktikleri.}
\label{tab:mathlib-taktik}
\end{table}

Bölüm~7'de Gauss formülünün son adımını, dağılma ve değişme teoremlerini
tek tek uygulayarak elle yapmıştık. \leantac{ring} bu işi tek adımda
yapar:

\leandosya[Gauss, \texttt{ring} ile]{Bolum08}{gauss_ring}

\noindent \Mathlib{}, sonlu toplamlar için matematikteki $\sum$
gösterimini de sunar. Aynı formül, kitaplarda yazıldığı biçimiyle:

\leandosya[Gauss, toplam gösterimiyle]{Bolum08}{gauss_finset}

\subsection{Teorem bulmak}

\Mathlib{}'de aradığınız teoremin büyük olasılıkla zaten var olduğunu
varsayın. Onu bulmanın üç yolu vardır:

\begin{enumerate}[adimlar]
  \item \textbf{Adını tahmin etmek.} \Mathlib{} katı bir adlandırma
        düzeni izler (\cref{tab:adlandirma}). ``Çarpma değişmelidir''
        teoremi \lc{mul\_comm}, ``bir kümenin kesişimine üyelik'' teoremi
        \lc{Set.mem\_inter\_iff}'tir.
  \item \textbf{\lc{exact?} ile aramak.} Hedefin yerine \lc{exact?}
        yazdığınızda \Lean{}, kütüphanede hedefi kapatan bir teorem arar
        ve bulursa önerir.
  \item \textbf{Çevrimiçi araçlar.} \Mathlib{}'in belgelendirme
        sayfaları ve teoremleri tiplerine göre arayan \emph{Loogle}
        arama motoru \citep{loogle}.
\end{enumerate}

\begin{leankod}[exact? ile arama (çıktı bir öneridir)]
example (a b : ℕ) : a + b = b + a := by
  exact?   -- Try this: exact Nat.add_comm a b
\end{leankod}

\begin{table}[htbp]
\centering
\renewcommand{\arraystretch}{1.4}
\rowcolors{2}{tablosatir}{white}
\begin{tabular}{>{\raggedright\arraybackslash}p{3.6cm}
                >{\raggedright\arraybackslash}p{4.6cm}
                >{\raggedright\arraybackslash}p{3.4cm}}
\rowcolor{tabloust}
\tablobaslik{Ad} & \tablobaslik{İfade} & \tablobaslik{Okunuşu} \\
\lc{mul\_comm}          & $a \cdot b = b \cdot a$          & çarpma, değişme \\
\lc{add\_assoc}         & $(a+b)+c = a+(b+c)$              & toplama, birleşme \\
\lc{mul\_add}           & $a(b+c) = ab + ac$               & çarpım, toplam üzerine \\
\lc{add\_zero}          & $a + 0 = a$                      & toplama, sıfır \\
\lc{le\_trans}          & $a \leq b \to b \leq c \to a \leq c$ & $\leq$, geçişme \\
\lc{Set.mem\_inter\_iff} & $x \in s \cap t \leftrightarrow x \in s \wedge x \in t$ & üyelik, kesişim, ancak ve ancak \\
\end{tabular}
\caption{\Mathlib{}'in adlandırma düzeninden örnekler. Ad, ifadenin
sembollerini soldan sağa okur; \lc{\_iff} bir denkliği, \lc{\_comm}
değişmeyi, \lc{\_assoc} birleşmeyi bildirir.}
\label{tab:adlandirma}
\end{table}

% ---------------------------------------------------------------------
\section{Kümeler yüklemdir}\label{sec:kume-yuklem}

\Mathlib{}'de bir küme, bir yüklemden başka bir şey değildir:

\begin{tanim}{Küme}{kume}
$\alpha$ tipinin bir \textbf{alt kümesi}, $\alpha$ üzerinde bir yüklemdir:
\[
  \mathsf{Set}\;\alpha \;:=\; \alpha \to \mathsf{Prop}.
\]
$x \in s$, ``$s(x)$ doğrudur'' demektir. $\{x \mid P(x)\}$ gösterimi,
$P$ yükleminin kendisidir.
\end{tanim}

\leandosya[Kümeler ve üyelik]{Bolum08}{kume_tanim}

\noindent İkinci örnekteki \lc{Iff.rfl}, ``$a \in \{x \mid P(x)\}$'' ile
``$P(a)$''nın \emph{tanım gereği} aynı olduğunu gösterir. \Cref{fig:kume-yuklem},
bu bakış açısını \cref{fig:yuklem}'deki yüklem resmiyle birleştiriyor.

\begin{figure}[htbp]
\centering
\begin{tikzpicture}[
    eleman/.style={circle, inner sep=2.2pt, fill=#1},
  ]
  \draw[anarenk, thick, fill=tanimbg, rounded corners=12pt] (0,0) rectangle (5.2,3.4);
  \node[font=\baslikfont\bfseries\small, text=anarenk, anchor=north west] at (0.15,3.35) {$\alpha = \mathbb{N}$};
  \draw[vurgu, very thick, fill=vurgu!12, rounded corners=10pt] (2.4,0.35) rectangle (4.9,2.45);
  \node[font=\baslikfont\bfseries\small, text=vurgu, anchor=south east] at (4.9,2.45) {$s = \{n \mid n \% 2 = 1\}$};
  \foreach \x/\y/\n in {2.9/1.9/1, 3.9/1.9/3, 2.9/0.9/5, 3.9/0.9/7}
    { \node[eleman=vurgu] at (\x,\y) {}; \node[font=\scriptsize, right=1pt] at (\x,\y) {$\n$}; }
  \foreach \x/\y/\n in {0.6/1.9/0, 1.5/1.9/2, 0.6/0.9/4, 1.5/0.9/6}
    { \node[eleman=black!35] at (\x,\y) {}; \node[font=\scriptsize, right=1pt] at (\x,\y) {$\n$}; }
  % karakteristik fonksiyon
  \node[font=\small, anchor=west, align=left] at (5.8,2.1)
    {$s(1) = (1 \% 2 = 1)$ \;\textcolor{bittiframe}{doğru}\\[2pt]
     $s(2) = (2 \% 2 = 1)$ \;\textcolor{hataframe}{yanlış}};
  \node[font=\baslikfont\scriptsize, text=black!60, anchor=west, align=left] at (5.8,0.8)
    {küme = her elemana bir önerme\\karşılık getiren fonksiyon};
\end{tikzpicture}
\caption{Bir küme, elemanlarını ``doğru'' önermelere götüren bir
yüklemdir.}
\label{fig:kume-yuklem}
\end{figure}

\begin{dikkat}[Dikkat: Tip kuramında kümeler homojendir]
Klasik küme kuramında her şey bir kümedir ve herhangi iki nesne
arasında $\in$ sorusu sorulabilir: ``$3 \in \pi$ mi?'' gibi anlamsız
sorular bile. Tip kuramında ise \lc{Set α} yalnızca \lc{α} tipindeki
elemanlardan oluşur; \lc{x ∈ s} yazabilmek için \lc{x : α} olmalıdır.
Bu kısıt, Russell paradoksunu (\cref{sec:berber}) daha baştan engeller:
``Kendini içermeyen kümelerin kümesi'' tanımlanamaz, çünkü bir küme
kendi tipinin elemanı değildir.
\end{dikkat}

% ---------------------------------------------------------------------
\section{Küme cebiri = önermeler mantığı}\label{sec:kume-cebiri}

Kümeler yüklem olduğuna göre, küme işlemleri de mantık bağlaçlarıdır
(\cref{tab:kume-mantik}).

\begin{table}[htbp]
\centering
\renewcommand{\arraystretch}{1.5}
\rowcolors{2}{tablosatir}{white}
\begin{tabular}{>{\centering\arraybackslash}p{1.6cm}
                >{\raggedright\arraybackslash}p{2.6cm}
                >{\raggedright\arraybackslash}p{4.6cm}
                >{\centering\arraybackslash}p{1.6cm}}
\rowcolor{tabloust}
\tablobaslik{Küme} & \tablobaslik{Adı} & \tablobaslik{Tanımı ($x \in \ldots$ ancak ve ancak)} & \tablobaslik{Mantık} \\
$s \cap t$       & kesişim      & $x \in s \wedge x \in t$ & $\wedge$ \\
$s \cup t$       & birleşim     & $x \in s \vee x \in t$   & $\vee$ \\
$s^{c}$          & tümleyen     & $\neg(x \in s)$          & $\neg$ \\
$s \setminus t$  & fark         & $x \in s \wedge \neg(x \in t)$ & $\wedge\neg$ \\
$s \subseteq t$  & alt küme     & her $x$ için $x \in s \to x \in t$ & $\to$ \\
$\emptyset$      & boş küme     & $\bot$                   & \lc{False} \\
\lc{univ}        & evrensel küme & $\top$                  & \lc{True} \\
\end{tabular}
\caption{Küme işlemleri ile mantık bağlaçları arasındaki sözlük.}
\label{tab:kume-mantik}
\end{table}

\begin{figure}[htbp]
\centering
\begin{tikzpicture}
  \def\A{(0,0) circle (0.75)}
  \def\B{(0.85,0) circle (0.75)}
  \foreach \i/\baslik in {0/{$s \cap t$}, 1/{$s \cup t$}, 2/{$s \setminus t$}, 3/{$s^{c}$}} {
    \begin{scope}[shift={({3.6*\i},0)}]
      \draw[black!20, rounded corners=4pt, fill=white] (-1.2,-1.1) rectangle (2.05,1.1);
      \ifcase\i
        \begin{scope}\clip \A; \fill[anarenk!45] \B; \end{scope}
      \or
        \fill[anarenk!45] \A; \fill[anarenk!45] \B;
      \or
        \begin{scope}\clip \A; \fill[anarenk!45] \A; \end{scope}
        \fill[white] \B;
      \or
        \fill[anarenk!45, rounded corners=4pt] (-1.2,-1.1) rectangle (2.05,1.1);
        \fill[white] \A;
      \fi
      \draw[thick, anarenk] \A;
      \draw[thick, vurgu] \B;
      \node[font=\scriptsize\bfseries, text=anarenk] at (-0.35,0.5) {$s$};
      \node[font=\scriptsize\bfseries, text=vurgu] at (1.2,0.5) {$t$};
      \node[font=\small] at (0.42,-1.4) {\baslik};
    \end{scope}
  }
\end{tikzpicture}
\caption{Dört temel küme işleminin Venn şemaları.}
\label{fig:venn}
\end{figure}

\begin{onemlifikir}[Önemli Fikir: Küme eşitlikleri mantık denklikleridir]
Bu sözlük, küme cebirindeki her yasanın önermeler mantığındaki bir
yasanın kılık değiştirmiş hâli olduğunu söyler. $s \cap (t \cup u) =
(s \cap t) \cup (s \cap u)$ dağılma yasası, Bölüm~2'deki
$P \wedge (Q \vee R) \leftrightarrow (P \wedge Q) \vee (P \wedge R)$
dağılma yasasıdır. George Boole'un 1854'te keşfettiği bu ortak yapıya
\emph{Boole cebiri} denir.
\end{onemlifikir}

\subsection{Alt küme ve küme eşitliği}

$s \subseteq t$, ``her $x$ için $x \in s$ ise $x \in t$'' demektir. Bir
alt küme ispatı, dolayısıyla \leantac{intro} ile başlar:

\leandosya[Alt kümenin geçişliliği]{Bolum08}{altkume}

\noindent İki kümenin eşitliği ise \emph{dışsallık} ilkesine dayanır:
Aynı elemanlara sahip iki küme eşittir. \leantac{ext x} taktiği,
\lc{s = t} hedefini \lc{x ∈ s ↔ x ∈ t} hedefine çevirir:

\leandosya[Kesişimin değişme özelliği]{Bolum08}{kesisim_degisme}

\begin{hedefgecisi}
\begin{hedefdurumu}[\texttt{ext x} öncesi]
s t : Set α
⊢ s ∩ t = t ∩ s
\end{hedefdurumu}
\sonra
\begin{hedefdurumu}[\texttt{ext x} sonrası]
s t : Set α
x : α
⊢ x ∈ s ∩ t ↔ x ∈ t ∩ s
\end{hedefdurumu}
\end{hedefgecisi}

\noindent \lc{x ∈ s ∩ t}, tanım gereği \lc{x ∈ s ∧ x ∈ t} olduğundan
\leantac{rintro ⟨hs, ht⟩} onu doğrudan parçalayabilir. Dağılma yasasının
ispatı, Bölüm~2'deki ispatın neredeyse harfi harfine aynısıdır:

\leandosya[Dağılma]{Bolum08}{dagilma}

\noindent Aynı ispatı \Mathlib{}'in araçlarıyla kısaltabiliriz:
\leantac{simp only} üyelik ifadelerini mantık bağlaçlarına çevirir,
\leantac{tauto} da ortaya çıkan totolojiyi kanıtlar:

\leandosya[Dağılma, otomasyonla]{Bolum08}{dagilma_simp}

\leandosya[De Morgan, kümeler için]{Bolum08}{tumleyen}

\noindent Bu ispat, Bölüm~3'teki \lc{demorgan\_veya} ispatının küme
diliyle yazılmış hâlidir. Kesişimin tümleyeni için De Morgan yasası ise,
Bölüm~3'teki ``zor yön'' nedeniyle klasik mantık gerektirir
(Alıştırma~\ref{alis:kume-demorgan}).

% ---------------------------------------------------------------------
\section{İlişkiler}\label{sec:iliskiler}

\begin{tanim}{İlişki}{iliski}
$\alpha$ üzerinde bir \textbf{ilişki}, iki argümanlı bir yüklemdir:
$R : \alpha \to \alpha \to \mathsf{Prop}$. $R\,a\,b$ ``$a$, $b$ ile
$R$ ilişkisindedir'' diye okunur.
\end{tanim}

\begin{table}[htbp]
\centering
\renewcommand{\arraystretch}{1.5}
\rowcolors{2}{tablosatir}{white}
\begin{tabular}{>{\raggedright\arraybackslash}p{2.6cm}
                >{\raggedright\arraybackslash}p{5.4cm}
                >{\raggedright\arraybackslash}p{3.6cm}}
\rowcolor{tabloust}
\tablobaslik{Özellik} & \tablobaslik{Tanım} & \tablobaslik{Örnek} \\
yansıma      & $\forall a,\; R\,a\,a$                                  & $=$, $\leq$, $\mid$ \\
simetri      & $\forall a\,b,\; R\,a\,b \to R\,b\,a$                   & $=$, ``kardeştir'' \\
ters simetri & $\forall a\,b,\; R\,a\,b \to R\,b\,a \to a = b$         & $\leq$, $\mid$ ($\mathbb{N}$'de) \\
geçişme      & $\forall a\,b\,c,\; R\,a\,b \to R\,b\,c \to R\,a\,c$    & $=$, $<$, $\leq$, $\mid$ \\
\end{tabular}
\caption{İlişkilerin temel özellikleri.}
\label{tab:iliski-ozellik}
\end{table}

Bu özelliklerin iki önemli bileşimi vardır:
\begin{itemize}
  \item \textbf{Denklik ilişkisi}: yansıma + simetri + geçişme.
        ``Aynı kalanı vermek'', ``aynı doğum gününe sahip olmak'' gibi
        ``bir bakıma aynı olmak'' ilişkileri.
  \item \textbf{Kısmi sıralama}: yansıma + ters simetri + geçişme.
        $\leq$, $\subseteq$ ve bölünebilirlik ($\mid$) gibi.
\end{itemize}

\subsection{Denklik ilişkileri}

``$n$'ye bölündüğünde aynı kalanı vermek'' bir denklik ilişkisidir.
\Lean{}'de bir denklik ilişkisi, üç özelliği alan olarak içeren bir
\lc{Equivalence} yapısıdır:

\leandosya[Mod $n$ denkliği]{Bolum08}{mod_denk}

\noindent Burada \leankw{where} ile bir yapının alanlarını tek tek
dolduruyoruz. Her alanın ispatı, eşitliğin ilgili özelliğinden gelir:
\lc{modDenk n a b}, tanım gereği \lc{a \% n = b \% n} eşitliği olduğu için
yansıma \lc{rfl}, simetri \lc{Eq.symm}, geçişme \lc{Eq.trans} ile
gösterilir.

Bir denklik ilişkisi, kümeyi \emph{denklik sınıflarına} ayırır: Birbirine
denk olan elemanlar aynı sınıftadır ve her eleman tam olarak bir sınıfa
aittir (\cref{fig:denklik-siniflari}).

\begin{figure}[htbp]
\centering
\begin{tikzpicture}[
    sinif/.style={draw=#1, fill=#1!10, very thick, rounded corners=10pt,
                  minimum width=3.4cm, minimum height=2.2cm},
    sayi/.style={font=\small},
  ]
  \node[sinif=anarenk]     (k0) at (0,0)   {};
  \node[sinif=ikincil]     (k1) at (4.0,0) {};
  \node[sinif=vurgu]       (k2) at (8.0,0) {};
  \node[font=\baslikfont\bfseries\small, text=anarenk, above=2pt of k0] {kalan 0};
  \node[font=\baslikfont\bfseries\small, text=ikincil, above=2pt of k1] {kalan 1};
  \node[font=\baslikfont\bfseries\small, text=vurgu, above=2pt of k2] {kalan 2};
  \node[sayi, align=center] at (k0) {$0 \quad 3 \quad 6$\\[3pt]$9 \quad 12 \;\cdots$};
  \node[sayi, align=center] at (k1) {$1 \quad 4 \quad 7$\\[3pt]$10 \quad 13 \;\cdots$};
  \node[sayi, align=center] at (k2) {$2 \quad 5 \quad 8$\\[3pt]$11 \quad 14 \;\cdots$};
\end{tikzpicture}
\caption{``3'e bölündüğünde aynı kalanı vermek'' ilişkisi, $\mathbb{N}$'yi
üç denklik sınıfına ayırır.}
\label{fig:denklik-siniflari}
\end{figure}

\subsection{Kısmi sıralamalar}

Bölünebilirlik, doğal sayılar üzerinde bir kısmi sıralamadır:

\leandosya[Bölünebilirlik]{Bolum08}{bolunme}

\noindent ``Kısmi'' sözcüğü, her iki elemanın karşılaştırılamayabileceğini
belirtir: Ne $2 \mid 3$ ne de $3 \mid 2$ doğrudur. Kısmi sıralamalar,
\emph{Hasse şemalarıyla} çizilir: $a \mid b$ ise $b$, $a$'nın üstünde
durur ve arada başka bir eleman yoksa ikisi bir çizgiyle birleştirilir
(\cref{fig:hasse}).

\begin{figure}[htbp]
\centering
\begin{tikzpicture}[
    d/.style={circle, draw=anarenk, fill=tanimbg, thick, minimum size=0.75cm, font=\small},
  ]
  \node[d] (1) at (0,0) {1};
  \node[d] (2) at (-1.6,1.4) {2};
  \node[d] (3) at (1.6,1.4) {3};
  \node[d] (4) at (-2.6,2.8) {4};
  \node[d] (6) at (0,2.8) {6};
  \node[d, draw=vurgu, fill=teorembg] (12) at (-1.3,4.2) {12};
  \foreach \a/\b in {1/2, 1/3, 2/4, 2/6, 3/6, 4/12, 6/12}
    \draw[thick, black!45] (\a) -- (\b);
  \node[font=\baslikfont\scriptsize, text=black!60, align=left, anchor=west] at (2.6,2.2)
    {12'nin bölenleri\\\emph{kısmi} sıralama:\\4 ile 6 karşılaştırılamaz};
\end{tikzpicture}
\caption{12'nin bölenlerinin bölünebilirlik sıralamasındaki Hasse şeması.}
\label{fig:hasse}
\end{figure}

% ---------------------------------------------------------------------
\section{Fonksiyonlar}\label{sec:fonksiyonlar}

\begin{tanim}{Birebir, örten, eşleme}{fonk-tur}
$f : \alpha \to \beta$ bir fonksiyon olsun.
\begin{itemize}
  \item $f$ \textbf{birebirdir} (İng.\ \emph{injective}):
        $\forall a\,b,\; f(a) = f(b) \to a = b$.
  \item $f$ \textbf{örtendir} (İng.\ \emph{surjective}):
        $\forall b,\; \exists a,\; f(a) = b$.
  \item $f$ \textbf{eşlemedir} (İng.\ \emph{bijective}): hem birebir hem örten.
\end{itemize}
\end{tanim}

\begin{figure}[htbp]
\centering
\begin{tikzpicture}[
    e/.style={circle, fill=#1, inner sep=1.8pt},
    ok/.style={-{Stealth[length=2mm]}, thick, black!50},
  ]
  \foreach \k/\baslik/\renk in {0/{birebir, örten değil}/anarenk, 1/{örten, birebir değil}/ikincil, 2/{eşleme}/vurgu} {
    \begin{scope}[shift={({4.4*\k},0)}]
      \draw[\renk, thick, fill=\renk!7, rounded corners=8pt] (-0.5,-1.6) rectangle (0.5,1.3);
      \draw[\renk, thick, fill=\renk!7, rounded corners=8pt] (1.9,-1.6) rectangle (2.9,1.3);
      \node[font=\small] at (0,1.6) {$\alpha$};
      \node[font=\small] at (2.4,1.6) {$\beta$};
      \node[font=\baslikfont\scriptsize\bfseries, text=\renk] at (1.2,-2.0) {\baslik};
    \end{scope}
  }
  % birebir, örten değil: 3 → 4
  \foreach \i/\y in {1/0.8, 2/0.0, 3/-0.8} \node[e=anarenk] (a\i) at (0,\y) {};
  \foreach \i/\y in {1/0.9, 2/0.3, 3/-0.3, 4/-0.9} \node[e=anarenk] (b\i) at (2.4,\y) {};
  \draw[ok] (a1) -- (b1); \draw[ok] (a2) -- (b2); \draw[ok] (a3) -- (b4);
  % örten, birebir değil: 4 → 3
  \foreach \i/\y in {1/0.9, 2/0.3, 3/-0.3, 4/-0.9} \node[e=ikincil] (c\i) at (4.4,\y) {};
  \foreach \i/\y in {1/0.8, 2/0.0, 3/-0.8} \node[e=ikincil] (d\i) at (6.8,\y) {};
  \draw[ok] (c1) -- (d1); \draw[ok] (c2) -- (d1); \draw[ok] (c3) -- (d2); \draw[ok] (c4) -- (d3);
  % eşleme: 3 → 3
  \foreach \i/\y in {1/0.8, 2/0.0, 3/-0.8} {
    \node[e=vurgu] (f\i) at (8.8,\y) {};
    \node[e=vurgu] (g\i) at (11.2,\y) {};
  }
  \draw[ok] (f1) -- (g2); \draw[ok] (f2) -- (g3); \draw[ok] (f3) -- (g1);
\end{tikzpicture}
\caption{Üç fonksiyon türü. Birebir fonksiyonda iki ok aynı yere gitmez;
örten fonksiyonda her hedef bir ok alır; eşlemede ikisi birden.}
\label{fig:fonksiyon-turleri}
\end{figure}

\Mathlib{}'de bu kavramlar \lc{Function.Injective}, \lc{Function.Surjective}
ve \lc{Function.Bijective} adını taşır. Tanımları, yukarıdaki
formüllerin harfi harfine aynısıdır; ispatlar da Bölüm~6'nın niceleyici
kurallarıyla yapılır:

\leandosya[İkiyle çarpma]{Bolum08}{birebir_ornek}

\noindent Birebirlik ispatı \leantac{intro a b h} ile başlar
($\forall$G, $\to$G); örten olmamanın ispatı ise varsayılan örtenliği
1 sayısına uygular ($\forall$E) ve tanığı \leantac{obtain} ile alır
($\exists$E).

\subsection{Bileşke}

\leandosya[Bileşke]{Bolum08}{bileske}

\noindent \lc{g ∘ f}, ``önce \lc{f}, sonra \lc{g}'' fonksiyonudur:
\lc{(g ∘ f) a = g (f a)}. Bu eşitlik tanımsal olduğu için \leantac{show}
ile hedefi açık biçimde yazabildik.

\subsection{Tersler}

\leandosya[Sol ters ve birebirlik]{Bolum08}{sol_ters}

\noindent İspattaki üç yeniden yazma adımı şu hesabı yapar:
$a = g(f(a)) = g(f(b)) = b$. Tersine, sağ tersi olan bir fonksiyon
örtendir (Alıştırma~\ref{alis:sag-ters}). Bu sonuçların \emph{tersleri}
ise ilginç bir biçimde asimetriktir (\cref{tab:tersler}).

\begin{table}[htbp]
\centering
\renewcommand{\arraystretch}{1.5}
\rowcolors{2}{tablosatir}{white}
\begin{tabular}{>{\raggedright\arraybackslash}p{5.2cm}
                >{\raggedright\arraybackslash}p{6.4cm}}
\rowcolor{tabloust}
\tablobaslik{Önerme} & \tablobaslik{Ne gerekir?} \\
sol tersi var $\Rightarrow$ birebir      & yalnızca mantık (yapıcı) \\
sağ tersi var $\Rightarrow$ örten        & yalnızca mantık (yapıcı) \\
birebir $\Rightarrow$ sol tersi var      & $\alpha$ boş olmamalı ve klasik mantık \\
örten $\Rightarrow$ sağ tersi var        & \textbf{seçim aksiyomu}: her $b$ için bir $a$ \emph{seçmek} \\
\end{tabular}
\caption{Tersler ve fonksiyon türleri.}
\label{tab:tersler}
\end{table}

\noindent Son satır özellikle önemlidir. $f$ örtense, her $b$ için
$f(a) = b$ olan \emph{bir} $a$ vardır; ama sağ tersi kurmak için her $b$
için \emph{belirli bir} $a$ seçmemiz gerekir. Sonsuz sayıda seçimi tek
seferde yapmak, seçim aksiyomunun tam olarak izin verdiği şeydir.
\Mathlib{}'de bu ters, \leankw{noncomputable} olarak işaretlenmiş
\lc{Function.surjInv} fonksiyonudur. Seçim aksiyomunu \cref{bol:aksiyomlar}'da
ayrıntılı olarak inceleyeceğiz.

% ---------------------------------------------------------------------
\section{Cantor teoremi}\label{sec:cantor}

\begin{teorem}{Cantor, 1891}{cantor}
Hiçbir $\alpha$ tipi için $\alpha$'dan $\mathsf{Set}\;\alpha$'ya örten bir
fonksiyon yoktur. Yani bir kümenin alt kümeleri, kümenin kendisinden
\emph{daha çoktur}.
\end{teorem}

\noindent Georg Cantor'un ünlü \emph{köşegen yöntemi}, berber
paradoksuyla aynı yapıya sahiptir. $f : \alpha \to \mathsf{Set}\;\alpha$
herhangi bir fonksiyon olsun ve şu kümeyi düşünelim:
\[
  D \;=\; \{x \mid x \notin f(x)\}.
\]
$D$, ``kendi görüntüsüne ait olmayan'' elemanların kümesidir. $f$ örten
olsaydı, $f(a) = D$ olan bir $a$ bulunurdu. O zaman $a \in D$ ancak ve
ancak $a \notin f(a) = D$ olurdu: Berber kendini tıraş eder mi?
(\cref{fig:kosegen})

\begin{figure}[htbp]
\centering
\begin{tikzpicture}[
    hucre/.style={draw=black!20, minimum size=0.75cm, font=\small},
    kosegen/.style={hucre, fill=vurgu!20, draw=vurgu, thick},
  ]
  \foreach \x/\v in {1/a_1, 2/a_2, 3/a_3, 4/a_4}
    \node[font=\small] at ({0.75*\x},0.75) {$\v$};
  \foreach \y/\v in {1/{f(a_1)}, 2/{f(a_2)}, 3/{f(a_3)}, 4/{f(a_4)}}
    \node[font=\small, anchor=east] at (0.25,{-0.75*(\y-1)}) {$\v$};
  \foreach \y/\satir in {1/{1,0,1,1}, 2/{0,0,1,0}, 3/{1,1,1,0}, 4/{0,1,0,0}} {
    \foreach \v [count=\x] in \satir {
      \ifnum\x=\y
        \node[kosegen] at ({0.75*\x},{-0.75*(\y-1)}) {$\v$};
      \else
        \node[hucre] at ({0.75*\x},{-0.75*(\y-1)}) {$\v$};
      \fi
    }
  }
  \node[font=\small, anchor=east] at (0.25,-3.25) {$D$};
  \foreach \v [count=\x] in {0,1,0,1}
    \node[hucre, fill=bittibg, draw=bittiframe, thick] at ({0.75*\x},-3.25) {$\v$};
  \draw[-{Stealth[length=2.4mm]}, thick, vurgu] (3.6,-1.1) to[bend left=30] (3.6,-3.1);
  \node[font=\baslikfont\scriptsize, text=vurgu, align=left, anchor=west] at (3.9,-2.0)
       {köşegeni\\tersine çevir};
  \node[font=\baslikfont\scriptsize, text=black!60, align=left, anchor=west] at (5.6,-0.6)
       {1: $a_j \in f(a_i)$\\0: $a_j \notin f(a_i)$};
  \node[font=\baslikfont\scriptsize, text=bittiframe, align=left, anchor=west] at (5.6,-3.25)
       {$D$, her satırdan en az\\bir yerde (köşegende) farklıdır};
\end{tikzpicture}
\caption{Cantor'un köşegen yöntemi. $D$ kümesi, $i$. satırla $i$.
sütunda farklı olacak biçimde kurulur; bu yüzden hiçbir satıra eşit
değildir.}
\label{fig:kosegen}
\end{figure}

\leandosya[Cantor teoremi]{Bolum08}{cantor}

\begin{tarihce}[Tarihçe: Sonsuzluğun büyüklükleri]
Cantor 1891'de yayımladığı kısa bir makalede köşegen yöntemini tanıttı
ve bir kümenin alt kümelerinin her zaman kümeden ``daha büyük'' olduğunu
gösterdi \citep{cantor1891}. Bu sonuç, farklı büyüklükte sonsuzlukların
var olduğu anlamına geliyordu ve döneminde büyük tartışmalara yol açtı.
Aynı köşegen fikri sonraki yarım yüzyılda Russell paradoksunda (1901),
Gödel'in eksiklik teoreminde (1931) ve Turing'in durma probleminde (1936)
yeniden karşımıza çıkacaktı.
\end{tarihce}

% ---------------------------------------------------------------------
\section{Taktik özeti}

\begin{center}
\renewcommand{\arraystretch}{1.45}
\rowcolors{2}{tablosatir}{white}
\begin{tabularx}{\linewidth}{>{\raggedright\arraybackslash}p{4.4cm} X}
\rowcolor{tabloust}
\tablobaslik{Taktik} & \tablobaslik{Ne yapar?} \\
\leantac{ext x}                    & Küme (ya da fonksiyon) eşitliğini elemanlara indirger. \\
\leantac{simp only [l₁, l₂]}       & Yalnızca verilen eşitliklerle sadeleştirir. \\
\leantac{tauto}                    & Önermeler mantığı totolojilerini kanıtlar (klasik). \\
\leantac{ring}                     & Değişmeli halkalarda polinom eşitlikleri. \\
\leantac{linarith}                 & Doğrusal eşitsizlikler. \\
\leantac{norm\_num}                & Sayısal hesaplar, asallık. \\
\leantac{positivity}               & Pozitiflik ve negatif olmama. \\
\lc{exact?}, \lc{apply?}           & Kütüphanede teorem arar. \\
\end{tabularx}
\end{center}

% ---------------------------------------------------------------------
\begin{arastirmanotu}
\textbf{Köşegen argümanlarının ortak yapısı.} F.\,William Lawvere 1969'da,
Cantor, Russell, Gödel ve Turing'in köşegen argümanlarının hepsinin
kategori kuramındaki tek bir sabit nokta teoreminin örnekleri olduğunu
gösterdi \citep{lawvere1969}. Bu sonucun matematik altyapısı ağır
olmayan, okunabilir bir sunumu için \citet{yanofsky2003}'ya bakılabilir.
Lawvere'in teoremi \Lean{}'de formalize edilmeye elverişli, kısa ve
öğretici bir projedir.

\textbf{Mathlib ve kütüphane tasarımı.} \Mathlib{}'in yapısı, adlandırma
düzeni ve ``tip sınıfı'' hiyerarşisi \citet{mathlib2020}'de anlatılır.
Büyük bir formal kütüphanede kavramların hangi genellikte tanımlanacağı,
hangi teoremlerin \leantac{simp} lemması yapılacağı gibi sorular aktif
tartışma konularıdır; bu konulara \ciltiki{9}'da döneceğiz.

\textbf{İleri okuma.} \Mathlib{}'le matematik yapmaya kapsamlı bir giriş
için \citet{mil}; kümeler ve fonksiyonlar bölümü bu bölümün doğal
devamıdır.
\end{arastirmanotu}

\begin{bolumozeti}
\begin{itemize}
  \item \Mathlib{}, \Lean{}'in büyük matematik kütüphanesidir. Teoremler
        katı bir adlandırma düzeni izler; \lc{exact?} ve Loogle arama
        araçlarıdır.
  \item \leantac{ring}, \leantac{linarith}, \leantac{norm\_num},
        \leantac{positivity}, \leantac{tauto} rutin işleri otomatikleştirir.
  \item Bir küme, bir yüklemdir: \lc{Set α := α → Prop}. Küme işlemleri
        mantık bağlaçlarıdır; küme cebiri önermeler mantığıdır.
  \item Küme eşitliği \leantac{ext} ile elemanlara indirgenir.
  \item Denklik ilişkileri sınıflara ayırır; kısmi sıralamalar Hasse
        şemalarıyla çizilir.
  \item Sol ters birebirliği, sağ ters örtenliği verir. Örten bir
        fonksiyonun sağ tersini kurmak seçim aksiyomu ister.
  \item Cantor teoremi: Hiçbir küme kendi alt kümeleri kadar büyük
        değildir. Köşegen yöntemi, berber paradoksunun genellemesidir.
\end{itemize}
\end{bolumozeti}

\begin{alistirmalar}
\alistirma{1} \lc{s ∩ s = s}. (\lc{a8\_1})

\alistirma{1} \lc{s ⊆ s ∪ t}. (\lc{a8\_2})

\alistirma{2} Birleşimin kesişim üzerine dağılması:
\lc{s ∪ (t ∩ u) = (s ∪ t) ∩ (s ∪ u)}. Önce elle, sonra \leantac{simp only}
ve \leantac{tauto} ile ispatlayın. (\lc{a8\_3})

\alistirma{2}\label{alis:kume-demorgan} (Klasik) Kesişimin tümleyeni:
\lc{(s ∩ t)ᶜ = sᶜ ∪ tᶜ}. Hangi yön klasik mantık gerektirir? Bölüm~3 ile
ilişkisini açıklayın. (\lc{a8\_4})

\alistirma{2} \lc{s ⊆ t → s ∩ u ⊆ t ∩ u}. (\lc{a8\_5})

\alistirma{2} Bölünebilirliğin geçişliliğini, \lc{dvd\_trans} kullanmadan
tanımdan ispatlayın. \lc{a ∣ b}, \lc{∃ c, b = a * c} demektir.
(\lc{a8\_6})

\alistirma{2} \lc{g ∘ f} örtense \lc{g} örtendir. (\lc{a8\_7})

\alistirma{2} \lc{g ∘ f} birebirse \lc{f} birebirdir. (\lc{a8\_8})

\alistirma{3}\label{alis:sag-ters} Sağ tersi olan fonksiyon örtendir:
\lc{(∀ b, f (k b) = b) → Surjective f}. (\lc{a8\_9})

\alistirma{3} Tam sayılar üzerinde ``farkı 3'ün katı olmak''
ilişkisinin bir denklik ilişkisi olduğunu ispatlayın. Kaç denklik sınıfı
vardır? (\lc{a8\_10})

\alistirma{3} Köşegen yöntemi: \lc{ℕ}'den \lc{ℕ → ℕ}'ye örten fonksiyon
yoktur. İpucu: \lc{fun n => F n n + 1} fonksiyonunu düşünün. Bu sonuç,
``bütün programlar sayılabilir ama bütün fonksiyonlar sayılamaz''
gözleminin özüdür. (\lc{a8\_11})

\alistirma{2} (Kâğıt-kalem) $\{1, 2, 3, 4, 6, 8, 12, 24\}$ kümesi için
bölünebilirlik sıralamasının Hasse şemasını çizin.

\alistirma{3} (Tartışma) \Cref{tab:tersler}'ün üçüncü satırında
``$\alpha$ boş olmamalı'' koşulu neden gerekir? $\alpha$ boşsa ve
$\beta$ boş değilse ne olur?
\end{alistirmalar}
